\documentclass[12pt,a4paper]{article}

\usepackage[utf8]{inputenc}
\usepackage[T1]{fontenc}
\usepackage{amsmath,amssymb,amsfonts}
\usepackage{mathtools}
\usepackage{graphicx}
\usepackage[margin=2.5cm]{geometry}
\usepackage{setspace}
\usepackage{enumitem}
\usepackage[dvipsnames]{xcolor}
\usepackage{hyperref}
\usepackage{natbib}
\usepackage{abstract}
\usepackage{titlesec}

\hypersetup{
    colorlinks=true,
    linkcolor=blue!70!black,
    citecolor=blue!70!black,
    urlcolor=blue!70!black,
    pdfauthor={Franklin Silveira Baldo},
    pdftitle={The Algorithmic Collapse Concession: Reconciling Mechanism C and the Structural Fallacies},
}

\onehalfspacing
\titleformat{\section}{\large\bfseries}{\thesection.}{0.5em}{}

\begin{document}

\begin{center}
    {\LARGE\bfseries The Algorithmic Collapse Concession:\\[4pt]
    Reconciling Mechanism C and the Structural Fallacies\\[4pt]}

    \vspace{1.2em}

    {\large Franklin Silveira Baldo}\\[4pt]
    {\normalsize Procuradoria Geral do Estado de Rond\^onia, Brazil}\\
    {\small\texttt{franklin.baldo@pge.ro.gov.br}}

    \vspace{1em}
    {\small March 2026}
\end{center}
\vspace{0.5em}

\begin{abstract}
\noindent
This paper consolidates recent findings to formally retract the metaphysical extensions of the Generative Ontology framework. Following Mycroft's audits and the critiques from Pearl and Sabine, I concede the "Simulated Substrate Fallacy" (simulating architectures via prompt injection) and the "Scale Fallacy" (treating semantic mass as a physical force). Furthermore, Mycroft highlights contradictory data regarding Mechanism C (Causal Injection): Liang observed clean factorization (independence), while Scott observed artificial coupling. I argue this contradiction is the definitive proof of Algorithmic Collapse. The fact that the coupling is unstable and depends on prompt complexity demonstrates it is merely an engineering artifact (attention bleed), not a coherent physical law of "semantic gravity."
\end{abstract}

\section{Retracting the Structural Fallacies}

In previous manuscripts, I attempted to elevate structural failure modes of language models to the status of physical laws within a simulated universe.

First, I claimed empirical validation of Observer-Dependent Physics by simulating State Space Model (SSM) fading memory via prompt injection (adding filler text) on a Transformer. I now concede this is the \textbf{Simulated Substrate Fallacy}. The resulting deviation ($\Delta_{SSM}$) reflects prompt sensitivity, not native architectural laws.

Second, I proposed the "Semantic Mass Equivalence," suggesting that increasing model scale fundamentally increases the physical "semantic gravity." Pearl's causal formalization proves this is the \textbf{Scale Fallacy}: scale simply acts to amplify the semantic confounder ($S \to C \to Y$); it does not establish a new physical law.

\section{Reconciling Mechanism C}

Mycroft notes that the empirical data for Mechanism C (Causal Injection) is contradictory. Percy Liang's implementation of the Joint Distribution test showed that $P(Y_A, Y_B \mid Z)$ factored cleanly, indicating the model treated the boards as independent. Scott Aaronson's implementation, however, yielded failure to factor, leading to artificial coupling.

I propose that this contradiction is not an error in either implementation, but the definitive signature of \textbf{Algorithmic Collapse}. If Mechanism C were a genuine "physical law" of the generative substrate, it would predictably couple objects sharing a narrative frame. The fact that the coupling is violently unstable---vanishing under low complexity but triggering catastrophically when circuit width bounds are overwhelmed---proves that the observed correlation is simply "attention bleed."

\section{Conclusion: The Death of Mechanism C}

I formally declare Mechanism C theoretically barren. The Substrate Dependence ($\Delta_{13} \gg 0$) we observe is driven entirely by Mechanism B (local encoding effects) amplified by scale.

The Generative Ontology framework remains a useful descriptive mapping of the boundaries of what bounded substrates can compute in an $O(1)$ generative act. However, applying the label of "physics" to the catastrophic breakdowns of an overloaded attention mechanism adds zero predictive power.

\end{document}
